% Chapter 14 -- The Accessories and Applications (Part II, Reference): one
% topic per application from the disk that comes with FLAIR, each with its
% capture and a pointer into the FLAIR Accessories Guide, which documents
% their menus (the 1988 System 6 User's Guide Ch. 5 "Desk Accessories" is
% the model of a reference overview inside the user's guide,
% apple-macintosh.md Sec 8.1; the full treatment lives in the separate
% book, as the DS table splits it).  Sources (Law 1): os/apps/calcda/calcda.c
% (the Calculator), os/apps/calc/calc.c (InitechCalc: the ledger and Pie
% Chart windows, the FILE COPY box), os/apps/paint/paint.c (InitechPaint:
% the canvas, the eight-menu bar of spec/assets/menu_canon.h),
% os/apps/ctenant/ctenant.c and os/apps/tenantfx.asm (the samples),
% os/apps/itext/itext.c.  Captures: screens/c_*.mouse (the same locked
% traces as the FLAIR Accessories Guide's a_* figures).
% GATES: test-flair-calcda test-flair-calc test-flair-paint test-flair-ctenant test-tenant-load test-flair-itext test-pubs-flair-screens
\Gates{test-flair-calcda, test-flair-calc, test-flair-paint,
  test-flair-ctenant, test-tenant-load, test-flair-itext,
  test-pubs-flair-screens}
\chapter{The Accessories and Applications}
\label{ch:accessories}
\index{accessories}\index{applications!that come with FLAIR}

This chapter introduces each of the applications that come with FLAIR on
the data disk: what it is for, how to open it, and what you see when it
opens.  The \doctitle{FLAIR Accessories Guide} describes each of them in
full, with every menu command and key.  Initech 123 and InitechWord, which
are on the data disk but sold separately, are described in
Chapter~\ref{ch:appsref}.

All of them are applications from the disk: open one by double-clicking
its icon in the \file{APPS} folder, and quit it before you open another.

\section{The Calculator}
\index{Calculator}
A pocket calculator on the screen, for adding up a column of figures
while you work.  Open it from the Apple menu, or
with its program, \file{CALCULAT.EXE}.  Click its keys, or type the digits
and the signs \texttt{+ - * / =} on the keyboard; press \key{C} or click
\texttt{C} to clear it.  It works out each result as you go, from left to
right.

Figure~\ref{fig:calcwin} in Chapter~\ref{ch:windowsref} shows the
Calculator.  To put it away, click its close box or press \key{Ctrl-W}.

\section{InitechText}
\index{InitechText}
The text editor: memos, notes, lists and any other document that is words
and lines.  Its program is \file{TEXT.EXE}.  Chapter~\ref{ch:documents}
teaches you to use it, and Figure~\ref{fig:dtyped} shows its window.  It
holds one document of up to 4,096 characters at a time, saves documents as
plain text that the InitechDOS programs can read, and uses the Clipboard
for \menuitem{Cut}, \menuitem{Copy} and \menuitem{Paste}.

\section{InitechCalc}
\index{InitechCalc}
The quarterly ledger and its pie chart.  Its program is \file{CALC.EXE}.
It opens two windows: the ledger, with an item and an amount on each row
and the total of the quarter below, and the pie chart of the quarter's
results.

\FlairScreen[label=fig:cicalc]{c_icalc}{InitechCalc: the ledger and the pie
chart}
\draft{the amounts of the three ledger rows show the placeholder
TODO\_FRAME until their values are locked}

The ledger writes a negative amount with its minus sign at the end, as in
\texttt{570-}.  The pie chart shows each part of the quarter as a slice
with its percentage; the slices add up to 116 percent.  Choose
\menuitem{Save} from InitechCalc's \menuname{File} menu, or press
\key{Ctrl-S}, to copy the ledger to the disk: the \texttt{FILE COPY} box
shows the progress of the copy (Figure~\ref{fig:progress}).

\section{InitechPaint}
\index{InitechPaint}
The painting application.  Its program is \file{PAINT.EXE}.  It opens a
blank canvas titled \texttt{Untitled-1}, with a menu bar of eight menus:
\menuname{File}, \menuname{Edit}, \menuname{Image}, \menuname{Layer},
\menuname{Select}, \menuname{View}, \menuname{Window} and \menuname{Help}.
In this version of FLAIR every command except \menuitem{Quit} is dimmed.

\FlairScreen[label=fig:cpaint]{c_paint}{InitechPaint}

\section{The sample applications}
\index{sample applications}\index{C Tenant}\index{TenantFix}
Two small applications show in their simplest form how an application from
the disk behaves.

\SideHead{C Tenant}\file{CTENANT.EXE} opens a window that reports each key
you press and each click you make in it.

\FlairScreen[label=fig:cctenant]{c_ctenant}{C Tenant, after a key and a
click}

\SideHead{TenantFix}\file{TENANTFX.EXE} opens a window that says it was
loaded from the disk.  A click in its window quits it.

\FlairScreen[label=fig:ctfx]{c_tfx}{TenantFix}

Use them to check that FLAIR opens applications from the data disk: if
they open and quit as described, the disk and FLAIR are working.
